\section{背景と課題}

\subsection{旧Python/Tkinter版の課題}

旧版（\texttt{font\_to\_binary\_GUI.py} / \texttt{font\_to\_binary\_CUI.py}）は
Python + Tkinter で実装された小規模ユーティリティであった。
一定の機能は果たしていたが、以下の課題があった。

\begin{itemize}[leftmargin=*,itemsep=2pt]
  \item \textbf{配布の煩雑さ}: Python と関連ライブラリ（Pillow, tkinter）を
        利用者がインストールする必要があった。
  \item \textbf{UI/UXの制約}: Tkinter のネイティブウィジェットでは
        近年の Web アプリのような細やかな表現ができず、
        レイアウトやテーマ切替も貧弱であった。
  \item \textbf{出力形式の限定}: C配列と一部の表記のみ対応しており、
        言語切替、ビットパック、MSB/LSB、0/1反転などの
        細かな設定ができなかった。
  \item \textbf{ピクセル編集不可}: 生成後の手動補正ができず、
        細かな調整をしたい場合は結果をコピーして手で書き換える必要があった。
  \item \textbf{共有性}: 設定を他人に渡す手段がなく、
        再現実験や共同作業に不向きであった。
\end{itemize}

\subsection{組み込み現場での典型的な需要}

マイコン用ディスプレイ（OLED/LCD/LED マトリクス）への
ドットフォント描画では、以下が繰り返し要求される。

\begin{itemize}[leftmargin=*,itemsep=2pt]
  \item ターゲットマイコン・表示ドライバごとに異なる
        バイト並び（行方向/列方向、MSB/LSB）への対応。
  \item FLASH 節約のためのビットパック配列出力（1bit/pixel）。
  \item \texttt{PROGMEM}, \texttt{const}, \texttt{static} 等の
        言語・処理系依存修飾子の付与。
  \item デザインを確認するためのプレビューと、
        生成結果の部分修正。
\end{itemize}

これらを毎回手作業で行うのは非効率であり、
\textbf{「設定をGUIで触り、結果をコピーするだけで済むツール」}
が常に求められていた。

\subsection{Webアプリ化の動機}

これらの課題に対し、
\textbf{クロスプラットフォームで即起動でき、機能拡張が容易で、
URLで設定を共有できる}という要求を満たすため、
Web アプリケーションとしての再実装を選択した。

\begin{itemize}[leftmargin=*,itemsep=2pt]
  \item 利用者はブラウザを開くだけで利用可能（インストール不要）。
  \item GitHub Pages による無料ホスティングで配布コスト \texttt{0}。
  \item React + TypeScript により型安全に機能拡張が可能。
  \item Canvas API + FontFace API によりフォントラスタライズを
        サーバレスで実行可能。
\end{itemize}

\newpage
